% Documento de detalles y notas de estudio que acompaña a la tesis; no forma parte del manuscrito.
% Se compila desde docs/ para que las referencias a la tesis resuelvan contra main.aux:
%   cd docs && latexmk -pdf -interaction=nonstopmode -outdir=<dir> chapters/_borradores/detalles.tex
% (compilar antes main.tex). Las tablas de calculos/ las genera
% benchmarks/np2_sequence/tesis/detalles_calculos.jl; no se editan a mano.
\documentclass[11pt,a4paper]{article}
\usepackage[T1]{fontenc}
\usepackage[spanish,es-noquoting,es-tabla]{babel}
\usepackage{lmodern}
\usepackage[a4paper,margin=2.4cm]{geometry}
\usepackage{amsmath,amssymb,bm}
\usepackage{physics}
\usepackage{booktabs,multirow,array}
\usepackage{siunitx}
\usepackage{graphicx}
\sisetup{group-digits=integer}
\spanishdecimal{.}
\usepackage{xcolor}
\usepackage{csquotes}
\usepackage{xr-hyper}
\usepackage[colorlinks=true,linkcolor=blue!50!black,citecolor=red!60!black]{hyperref}
\usepackage[backend=biber,style=apa,sorting=nyt]{biblatex}
\addbibresource{bib/bibliografia.bib}
\externaldocument{main}
\setlength{\parskip}{0.4em}

\title{Detalles adicionales y notas de estudio\\[0.3em]
\large Complemento de la tesis; no forma parte del manuscrito}
\author{}
\date{Octubre de 2026}

\begin{document}
\maketitle
\vspace{-2em}

Este documento reúne lo que se retiró del manuscrito o nunca cupo en él: desarrollos que la tesis
cita solo como resultado, los cálculos intermedios que hay detrás de algunos números del
Capítulo~\ref{ch:resultados} y notas de estudio para el examen. Los números de sección, ecuación y
tabla que aparecen aquí son los de la tesis. Las tablas de cálculos las genera el script
\texttt{detalles\_calculos.jl}, en \texttt{benchmarks/np2\_sequence/tesis/}, a partir de los mismos resultados
registrados que la etapa~3, con sus mismas funciones, de modo que ningún número de este documento
está escrito a mano.

\tableofcontents

%======================================================================
\section{Acoplamiento intermedio y teoría de perturbaciones}
\label{det:perturbaciones}

La Sección~\ref{sec:acoplamiento_ls} de la tesis dice que la mezcla entre términos puede
estimarse con teoría de perturbaciones mientras $\zeta_{np}$ es pequeño, y que por eso el
espín-órbita se diagonaliza de forma directa. Aquí está el argumento completo.

\subsection{El bloque de dos niveles}

Cada bloque de $J$ con dos estados (Ecuaciones~\eqref{eq:bp_j0} y \eqref{eq:bp_j2}) tiene la forma
\begin{equation}
  \bm{H} = \begin{pmatrix} a & x \\ x & b \end{pmatrix}, \qquad \delta = b - a > 0,
\end{equation}
donde $a$ es el elemento diagonal del triplete, $b$ el del singlete y $x$ el elemento de
espín-órbita que los conecta. La teoría de perturbaciones de Rayleigh-Schrödinger separa
$\bm{H} = \bm{H}_0 + \bm{V}$ con $\bm{H}_0$ diagonal y $\bm{V}$ el elemento fuera de la diagonal.
Como $\bm{V}$ no tiene elementos diagonales, la corrección de primer orden a la energía se anula y
las primeras correcciones son
\begin{equation}
  E^{(2)} = -\frac{x^2}{\delta}, \qquad
  \ket{\psi} \simeq \ket{a} - \frac{x}{\delta}\ket{b}
  \quad\Longrightarrow\quad w^{(1)} = \frac{x^2}{\delta^2},
\end{equation}
donde $w$ es el peso del singlete en el nivel del triplete. Es la estimación de la
Ecuación~\eqref{eq:mezcla_perturbativa}.

\subsection{La solución exacta y el radio de convergencia}

El bloque se diagonaliza en forma cerrada. Con $z = 2x/\delta$,
\begin{equation}
  E_- = a - \frac{\delta}{2}\qty(\sqrt{1+z^2} - 1), \qquad
  w = \frac{1}{2}\qty(1 - \frac{1}{\sqrt{1+z^2}}).
\end{equation}
La serie perturbativa no es otra cosa que el desarrollo de Taylor de $\sqrt{1+z^2}$ en potencias
de $z$:
\begin{equation}
  E_- = a - \frac{x^2}{\delta} + \frac{x^4}{\delta^3} - 2\frac{x^6}{\delta^5} + \cdots, \qquad
  w = \frac{z^2}{4} - \frac{3z^4}{16} + \cdots .
\end{equation}
La función $\sqrt{1+z^2}$ tiene puntos de ramificación en $z = \pm i$, así que su serie de Taylor
converge solo si $|z| < 1$, es decir, si $|x| < \delta/2$. Esa es la versión precisa de \emph{la
serie perturbativa deja de ser fiable}. Para $|z| \ll 1$ basta el primer término. Cuando $|z|$ se
acerca a uno, los términos decrecen lentamente y truncar la serie cambia el resultado. En el límite
$jj$, donde la separación electrostática desaparece y $\delta \to 0$, se tiene $|z| \to \infty$ y
la serie diverge. La diagonalización exacta vale en todos los casos, y por eso la tesis la usa.

\subsection{Los números de la secuencia}

La Tabla~\ref{tab:det_mezcla} evalúa ambos caminos con los mismos elementos de matriz que
ensambla el motor, incluido el factor de los coeficientes CI de la Sección~\ref{sec:ci_method}.

\begin{table}[htbp]
  \centering\small
  \input{chapters/_borradores/calculos/mezcla_perturbativa}
  \caption{Bloques $J=2$ y $J=0$ con las energías del CI y el $\zeta_{np}$ calculado (cm$^{-1}$).
  $w^{(1)}$ y $E^{(2)}$ son los resultados de primer y segundo orden; $w$ y $E_{\text{exacto}}$, los
  de la diagonalización. El $w$ del bloque $J=2$ es el peso de $^1D_2$ que da la
  Tabla~\ref{tab:lande} para el modelo.}
  \label{tab:det_mezcla}
\end{table}

En el carbono y el silicio, $z$ no pasa de 0.03 y el primer orden coincide con la solución exacta
en todas las cifras mostradas. En el germanio, $z \approx 0.18$ y el primer orden sobrestima el
peso en un par por ciento relativo. En el bloque $J=2$ del estaño, $z \approx 0.5$: la serie aún
converge, pero el primer orden da un peso de 6.5\,\% frente al 5.5\,\% exacto, un error relativo de
casi una quinta parte.

La Ecuación~\eqref{eq:mezcla_perturbativa} de la tesis usa como denominador la separación
$\Delta_{J=2}$ entre los niveles \emph{finales}, no la separación $\delta$ entre los elementos
diagonales. Como $\Delta = \delta\sqrt{1+z^2}$, ese cambio absorbe parte de los órdenes
superiores: en el estaño da un peso cercano al 5.2\,\% con los números de esta tabla, mucho más
próximo al exacto que el 6.5\,\% de primer orden estricto. Por eso la Sección~\ref{sec:lande}
puede decir que la estimación reproduce casi los pesos de la diagonalización.

%======================================================================
\section{Ejercicios sobre las matrices de Breit-Pauli}
\label{det:ejercicios_bp}

\subsection{Qué conserva el espín-órbita}

\paragraph{Enunciado.} Mostrar que $\hat{H}_{SO} = \zeta\,(\mathbf{l}_1\cdot\mathbf{s}_1 + \mathbf{l}_2\cdot\mathbf{s}_2)$
conmuta con $\hat{J}^2$ y no con $\hat{L}^2$ ni con $\hat{S}^2$.

\paragraph{Un electrón.} Con $\mathbf{j} = \mathbf{l} + \mathbf{s}$ se tiene
$\mathbf{l}\cdot\mathbf{s} = \tfrac{1}{2}(\mathbf{j}^2 - \mathbf{l}^2 - \mathbf{s}^2)$. Las tres piezas
conmutan con cada componente de $\mathbf{j}$: $\mathbf{j}^2$ porque es su invariante, y
$\mathbf{l}^2$, $\mathbf{s}^2$ porque cada una conmuta con su propio momento y con el del otro
espacio. Así que $[\mathbf{l}\cdot\mathbf{s}, \mathbf{j}] = 0$. Nótese que para un solo electrón
$\mathbf{l}\cdot\mathbf{s}$ \emph{sí} conmuta con $\mathbf{l}^2$ y $\mathbf{s}^2$: lo que el
espín-órbita no conserva es la \emph{dirección} de $\mathbf{l}$ y de $\mathbf{s}$, no su magnitud.

\paragraph{Dos electrones: $J$ se conserva.} Cada sumando $\mathbf{l}_i\cdot\mathbf{s}_i$ conmuta con
$\mathbf{j}_i$ por el punto anterior, y con $\mathbf{j}_k$ ($k\neq i$) porque actúan sobre electrones
distintos. Entonces $\hat{H}_{SO}$ conmuta con $\mathbf{J} = \mathbf{j}_1 + \mathbf{j}_2$ y con
$\mathbf{J}^2$.

\paragraph{Dos electrones: $L$ no se conserva.} Ahora $\mathbf{L}^2 = \mathbf{l}_1^2 + \mathbf{l}_2^2 +
2\,\mathbf{l}_1\cdot\mathbf{l}_2$, y el término cruzado es el que falla. Con
$[l_{1a}, l_{1b}] = i\epsilon_{abc}\,l_{1c}$ (en unidades de $\hbar$),
\begin{equation}
  [\mathbf{l}_1\cdot\mathbf{s}_1,\ \mathbf{l}_1\cdot\mathbf{l}_2]
  = \sum_{a,b} s_{1a}\,l_{2b}\,[l_{1a}, l_{1b}]
  = i \sum_{abc}\epsilon_{abc}\,s_{1a}\,l_{2b}\,l_{1c}
  = i\,(\mathbf{s}_1\times\mathbf{l}_2)\cdot\mathbf{l}_1 ,
\end{equation}
que no es el operador nulo. La forma más rápida de verlo, sin embargo, es con la
Tabla~\ref{tab:so_racah}: si $\hat{H}_{SO}$ conmutara con $\hat{L}^2$, sus elementos entre estados
de distinto $L$ serían cero, porque
$\mel{L}{[\hat{H}_{SO},\hat{L}^2]}{L'} = [L'(L'+1) - L(L+1)]\mel{L}{\hat{H}_{SO}}{L'}$. Pero
$\mel{^3P_0}{\hat{H}_{SO}}{^1S_0} = \mp\sqrt{2}\,\zeta \neq 0$ conecta $L=1$ con $L=0$ (y $S=1$
con $S=0$), así que no conmuta ni con $\hat{L}^2$ ni con $\hat{S}^2$. En la identidad~\eqref{eq:so_identidad}
el responsable es el segundo sumando, $(\mathbf{l}_1-\mathbf{l}_2)\cdot(\mathbf{s}_1-\mathbf{s}_2)$.

\subsection{El cálculo directo con símbolos \texorpdfstring{$6j$}{6j}, paso a paso}

La Sección~\ref{sec:so_racah} de la tesis da la fórmula~\eqref{eq:so_racah}. Para el elemento
$^3P_0$--$^1S_0$ ($L=1$, $S=1$, $J=0$; $L'=0$, $S'=0$) los factores son:
\begin{itemize}
\item La fase: $(-1)^{L'+S+J} = (-1)^{0+1+0} = -1$.
\item El $6j$ del acoplamiento: $\{1\,1\,0;\,0\,0\,1\}$. Con un cero, la fórmula general es
$\{a\,b\,0;\,d\,e\,f\} = (-1)^{a+e+f}\,\delta_{ab}\,\delta_{de}/\sqrt{(2a+1)(2d+1)}$, que da
$(-1)^{2}/\sqrt{3} = 1/\sqrt{3}$.
\item La parte orbital, con la Ecuación~\eqref{eq:reducido_un_electron} y $j=l=1$, $J_{12}=1$, $J'_{12}=0$:
fase $(-1)^{2+0+1} = -1$, raíz $\sqrt{3\cdot1}$, el $6j$ $\{1\,1\,1;\,0\,1\,1\} = -1/3$, y
$\sqrt{l(l+1)(2l+1)} = \sqrt{6}$. El producto es $(-1)\sqrt{3}\,(-1/3)\sqrt{6} = \sqrt{2}$.
\item La parte de espín, con $j=s=1/2$, $J_{12}=1$, $J'_{12}=0$: fase $(-1)^{1+0+1} = +1$, raíz
$\sqrt{3}$, el $6j$ $\{\tfrac12\,1\,\tfrac12;\,0\,\tfrac12\,1\} = 1/\sqrt{6}$, y $\sqrt{\tfrac12\cdot\tfrac32\cdot2} = \sqrt{3/2}$.
El producto es $\sqrt{3}/2$.
\end{itemize}
Todo junto, $2\zeta\,(-1)(1/\sqrt{3})(\sqrt{2})(\sqrt{3}/2) = -\sqrt{2}\,\zeta$. Los $6j$ con un cero
salen de la fórmula de arriba tras permutar columnas, que no cambia el valor; los demás
($\{1\,1\,1;\,1\,1\,1\} = 1/6$, $\{1\,1\,1;\,2\,1\,1\} = 1/6$) conviene buscarlos en una tabla o
calcularlos con \texttt{WignerSymbols.jl}. En el repositorio, el script de comprobación
\texttt{benchmarks/np2\_sequence/check\_so\_matrix\_elements.jl}\linebreak imprime los cinco
elementos de la Tabla~\ref{tab:so_racah} con otra convención de fases.

\subsection{Los mismos valores a partir del acoplamiento \texorpdfstring{$jj$}{jj}}

\paragraph{Enunciado.} Sin álgebra de Racah, obtener $|x|$ en los bloques $J=0$ y $J=2$ a partir
de la Sección 10-4 de \textcite{Cowan1981}.

\paragraph{Solución.} En la base $jj$, donde cada electrón tiene $j = l\pm\tfrac12$ definido,
$\hat{H}_{SO}$ es diagonal: como $\mathbf{l}\cdot\mathbf{s} = \tfrac12[j(j+1)-l(l+1)-s(s+1)]$, un
electrón $p_{1/2}$ aporta $-\zeta$ y uno $p_{3/2}$, $+\tfrac12\zeta$. Cowan lo enuncia para electrones
no equivalentes, pero vale igual para $np^2$, porque $\hat{H}_{SO}$ es suma de operadores de un
cuerpo diagonales en los orbitales de $j$ definido; el principio de exclusión solo restringe los
$J$: $(p_{1/2})^2$ da $J=0$, $p_{1/2}p_{3/2}$ da $J=1,2$ y $(p_{3/2})^2$ da $J=0,2$. Los eigenvalores
de $\hat{H}_{SO}$ son entonces $\{-2\zeta, +\zeta\}$ en $J=0$ y $\{-\tfrac12\zeta, +\zeta\}$ en $J=2$.
En la base $LS$ el mismo operador tiene diagonal $(-\zeta, 0)$ y $(+\tfrac12\zeta, 0)$
(Sección~\ref{sec:matrices_bp}). Traza y determinante no dependen de la base: la traza se cumple
($-\zeta = -2\zeta+\zeta$ y $\tfrac12\zeta = -\tfrac12\zeta + \zeta$), y el determinante $-x^2$ debe
ser el producto de los eigenvalores, lo que da $x^2 = 2\zeta^2$ y $x^2 = \zeta^2/2$. El signo de $x$
no queda fijado, y eso es justamente la libertad de fase.

%======================================================================
\section{Relación con la teoría de perturbaciones de Møller-Plesset}
\label{det:mp}

No conozco la pregunta exacta que hizo tu asesor. Lo que sigue son las dos conexiones que hay entre
este trabajo y la teoría de Møller-Plesset, que \textcite[Capítulo 6]{Szabo1989} desarrolla.

\paragraph{Qué corrige MP2.} La teoría de Møller-Plesset es la de Rayleigh-Schrödinger con
$\hat{H}_0$ igual a la suma de los operadores de Fock. Con esa partición, $E^{(0)} + E^{(1)}$ es
exactamente la energía de Hartree-Fock, y la primera corrección es la de segundo orden, que
recupera correlación a partir de las excitaciones dobles:
\begin{equation}
  E^{(2)}_{\text{MP}} = \sum_{i<j}\sum_{a<b}
  \frac{\abs{\langle ij|ab\rangle - \langle ij|ba\rangle}^2}{\epsilon_i + \epsilon_j - \epsilon_a - \epsilon_b},
\end{equation}
con $i, j$ ocupados, $a, b$ virtuales y los elementos de matriz de $1/r_{12}$.

\paragraph{Primera conexión: el CI de pareja.} El CI de la Sección~\ref{sec:ci_teoria}
diagonaliza el Hamiltoniano en el espacio de las funciones de pareja de valencia. Tratar esa misma
matriz con Rayleigh-Schrödinger desde la referencia $np^2$ da, a segundo orden,
$E^{(2)} = \sum_K \abs{H_{0K}}^2/(E_0 - E_K)$, que es la energía de correlación de la pareja de
valencia a segundo orden. Si los denominadores se toman de energías orbitales es la contribución
MP2 de esa pareja; si se toman de la diagonal de la matriz, es la partición de Epstein-Nesbet. La
diagonalización del CI de pareja contiene todos los órdenes dentro de ese espacio, del mismo modo
que la diagonalización del bloque de espín-órbita contiene toda la serie de la
Sección~\ref{det:perturbaciones}.

\paragraph{Segunda conexión: el cambio de base.} Diagonalizar un bloque es girar la base hasta
la de sus autovectores. En el bloque de dos niveles el ángulo exacto cumple
$\tan 2\theta = 2x/\delta$, y la teoría de perturbaciones lo aproxima a primer orden por
$\theta \simeq x/\delta$, que es el coeficiente de mezcla de la sección anterior. Es probablemente
el \emph{cambio de base} que recuerdas: la ecuación secular del bloque de espín-órbita es la misma
que se resuelve en la perturbación, pero sin truncar.

%======================================================================
\section{Los intervalos de Landé}
\label{det:lande}

\subsection{La regla}

Dentro de un término $LS$, cualquier operador de la forma $\sum_i \xi(r_i)\,\mathbf{l}_i\cdot\mathbf{s}_i$
es un producto escalar de un tensor de rango 1 en el espacio orbital por otro de rango 1 en el de
espín. Por el teorema de Wigner-Eckart, su valor esperado en $\ket{LSJM}$ depende de $J$ solo a
través de un símbolo $6j$ con un $1$ en la esquina, que es el mismo que el de $\mathbf{L}\cdot\mathbf{S}$.
Por eso, dentro del término, el operador equivale a $A\,\mathbf{L}\cdot\mathbf{S}$ con una
constante $A$ que depende del término pero no de $J$:
\begin{equation}
  E_J = E_{LS} + \frac{A}{2}\qty[J(J+1) - L(L+1) - S(S+1)]
  \quad\Longrightarrow\quad E_J - E_{J-1} = A\,J .
\end{equation}
Es la \textbf{regla de intervalos de Landé}: los intervalos crecen linealmente con el $J$ superior.

\subsection{El signo de \texorpdfstring{$A$}{A} y el triplete de \texorpdfstring{$p^2$}{p2}}

Para el término fundamental de una subcapa $l^w$, $A = +\zeta/2S$ si la subcapa está llena a
menos de la mitad y $A = -\zeta/2S$ si está llena a más de la mitad. En $p^2$ el $^3P$ tiene
$S=1$ y $A = \zeta_{np}/2$, el coeficiente del primer sumando de la Ecuación~\eqref{eq:so_identidad}:
el nivel de menor $J$ queda abajo y el multiplete es \emph{normal}. En $p^4$, como el oxígeno, el
signo se invierte y el $^3P_2$ es el fundamental: el multiplete es \emph{invertido}. Es la tercera
regla de Hund de la Sección~\ref{sec:np2_structure}.

Con $J = 0, 1, 2$, los intervalos del $^3P$ valen $A$ y $2A$, de modo que la razón
\begin{equation}
  R = \frac{E(^3P_2) - E(^3P_1)}{E(^3P_1) - E(^3P_0)} = 2
\end{equation}
no depende de $\zeta_{np}$. Por eso la Sección~\ref{sec:regimen_acoplamiento} la usa como
indicador: es insensible a la escala del espín-órbita y solo se aparta de 2 por mecanismos que
rompen la forma $A\,\mathbf{L}\cdot\mathbf{S}$.

\subsection{Qué rompe la regla}

\begin{enumerate}
\item \textbf{La mezcla con otros términos del mismo $J$.} En $p^2$ el $^3P_2$ se mezcla con el
$^1D_2$ y el $^3P_0$ con el $^1S_0$, mientras que el $^3P_1$ no tiene compañero. A segundo orden,
los dos niveles extremos bajan en $x^2/\delta$ (Tabla~\ref{tab:det_mezcla}) y el intermedio no se
mueve. Como el $^3P_2$ y el $^3P_0$ están en extremos opuestos del triplete, la bajada del
$^3P_2$ acorta el intervalo superior y la del $^3P_0$ alarga el inferior: los dos efectos bajan
$R$. Es el mecanismo del germanio y del estaño.
\item \textbf{La interacción espín-otra-órbita.} También es un producto escalar de rango 1 en el
espacio orbital por rango 1 en el de espín, de modo que dentro de un término solo cambia el valor
de $A$ y respeta la regla.
\item \textbf{La interacción espín-espín.} Es un producto de tensores de rango 2 en cada espacio.
Su dependencia en $J$ dentro del término la da un $6j$ con un 2 en la esquina, que no es la de
$\mathbf{L}\cdot\mathbf{S}$, así que rompe la regla aun sin mezcla. Es la candidata de la tesis
para el carbono, donde la mezcla es despreciable y $R$ medido vale 1.64.
\end{enumerate}

\begin{table}[htbp]
  \centering
  \input{tablas/acoplamiento}
  \caption{La Tabla~\ref{tab:acoplamiento} de la tesis: escalas del régimen de acoplamiento
  (cm$^{-1}$) y razón de intervalos del modelo y del NIST.}
\end{table}

El modelo da $R \simeq 2$ en el carbono porque su único operador de estructura fina es de un cuerpo
y la mezcla es despreciable ($z \approx 0.005$). En el germanio y el estaño reproduce la caída de
$R$ porque esa caída es la mezcla del punto~1. Lo que no reproduce es el 1.64 del carbono, que
exige el punto~3.

%======================================================================
\section{Extrapolación de las curvas de convergencia del CI}
\label{det:extrapolacion}

La Sección~\ref{sec:convergencia_ci} extrapola los singletes con los tres últimos tamaños de cada
curva. Aquí están los dos modelos, sus fórmulas y todos los números intermedios.

\subsection{Los dos modelos}

Sean $E_1, E_2, E_3$ las energías de un término en los tamaños $m_1 < m_2 < m_3$, y $d_1 = E_2 -
E_1$, $d_2 = E_3 - E_2$ los incrementos.

\paragraph{Geométrico.} Supone que cada incremento es el anterior por una razón constante
$r = d_2/d_1$. Si $0 < r < 1$, los incrementos que faltan son $d_2 r + d_2 r^2 + \cdots$, una serie
geométrica, y
\begin{equation}
  E_\infty^{\text{geo}} = E_3 + d_2\,\frac{r}{1-r}.
\end{equation}

\paragraph{De potencia.} Supone $E(m) = E_\infty + A\,m^{-p}$, con tres incógnitas que los tres
puntos fijan exactamente. El cociente de incrementos elimina $E_\infty$ y $A$,
\begin{equation}
  \frac{d_1}{d_2} = \frac{m_1^{-p} - m_2^{-p}}{m_2^{-p} - m_3^{-p}},
\end{equation}
una ecuación en $p$ que se resuelve por bisección. Después,
$A = d_2/(m_3^{-p} - m_2^{-p})$ y $E_\infty^{\text{pot}} = E_3 - A\,m_3^{-p}$.

Ninguno de los dos modelos se deduce de la teoría. Para la expansión en ondas parciales del helio
se conocen leyes de potencia, pero aquí $m$ cuenta funciones radiales por canal y no momentos
angulares. Por eso se usan los dos: cuando coinciden, la cola de convergencia está bien acotada;
cuando difieren, la diferencia mide la incertidumbre de la extrapolación.

Las Tablas~\ref{tab:det_extrap_datos} y \ref{tab:det_extrap_res} recogen todos los números de las
ocho curvas.

\begin{table}[htbp]
  \centering\small
  \input{chapters/_borradores/calculos/extrapolaciones_datos}
  \caption{Los tres últimos tamaños de cada curva de singletes (cm$^{-1}$, referidos al $^3P_0$),
  sus incrementos y la razón entre ellos.}
  \label{tab:det_extrap_datos}
\end{table}

\begin{table}[htbp]
  \centering\small
  \input{chapters/_borradores/calculos/extrapolaciones_resultados}
  \caption{Extrapolaciones geométrica y de potencia (cm$^{-1}$), con el error relativo frente al
  NIST entre paréntesis. La tercera columna es el error del último tamaño calculado.}
  \label{tab:det_extrap_res}
\end{table}

\subsection{El caso del \texorpdfstring{$^1D_2$}{1D2} del estaño}

Es el caso que la tesis declara indeterminado, y las tablas muestran por qué. Sus dos incrementos
casi no cambian: $d_1 = -28.1$ y $d_2 = -26.0$ cm$^{-1}$, con $r = 0.93$. El factor $r/(1-r)$ vale
entonces cerca de 12, de modo que la extrapolación geométrica multiplica el último incremento por
doce y baja el nivel unos 320 cm$^{-1}$, hasta 8388.8, un 2.6\,\% por debajo del NIST. Con $r$ tan
cerca de 1, un cambio pequeño en un incremento cambia mucho ese factor, y la extrapolación no es
fiable.

El modelo de potencia ni siquiera tiene solución, y por eso su fila lleva guiones. Cuando
$p \to 0$, el cociente $(m_1^{-p}-m_2^{-p})/(m_2^{-p}-m_3^{-p})$ tiende a
$\ln(m_2/m_1)/\ln(m_3/m_2) = \ln(4/3)/\ln(5/4) \approx 1.29$, y crece con $p$. El cociente medido,
$d_1/d_2 = 1.08$, queda por debajo de ese mínimo: los incrementos decrecen más despacio que
cualquier ley $m^{-p}$ con $p > 0$. Con el espacio de producción no se puede decir si el $^1D_2$
del estaño acabará por encima o por debajo del valor medido.

%======================================================================
\section{El ajuste de ley de potencias frente a \texorpdfstring{$Z$}{Z}}
\label{det:potencias}

La Sección~\ref{sec:regimen_acoplamiento} resume el crecimiento de $\zeta_{np}$ con exponentes
de una ley $y = A\,Z^p$. El ajuste es por mínimos cuadrados sobre los logaritmos: con
$x_i = \ln Z_i$ y $y_i = \ln Y_i$,
\begin{equation}
  p = \frac{\sum_i (x_i - \bar{x})(y_i - \bar{y})}{\sum_i (x_i - \bar{x})^2}, \qquad
  \ln A = \bar{y} - p\,\bar{x}, \qquad
  R^2 = 1 - \frac{\sum_i (y_i - \ln A - p\,x_i)^2}{\sum_i (y_i - \bar{y})^2}.
\end{equation}
El exponente local entre dos elementos consecutivos es
$\ln(Y_{i+1}/Y_i)/\ln(Z_{i+1}/Z_i)$.

\begin{table}[htbp]
  \centering\small
  \input{chapters/_borradores/calculos/ley_potencias_datos}
  \caption{Datos del ajuste: las tres $\zeta$ y $F^2$ en cm$^{-1}$, $\langle r^{-3}\rangle$ en
  $a_0^{-3}$.}
  \label{tab:det_pot_datos}
\end{table}

\begin{table}[htbp]
  \centering\small
  \input{chapters/_borradores/calculos/ley_potencias_ajuste}
  \caption{Exponente $p$, prefactor $A$ y $R^2$ del ajuste sobre los cuatro elementos, y
  exponentes locales entre elementos consecutivos.}
  \label{tab:det_pot_ajuste}
\end{table}

Tres observaciones sostienen lo que dice la tesis.

\begin{enumerate}
\item \textbf{La diferencia de una unidad es exacta.} Como
$\zeta_{\text{desnudo}} = (\alpha^2/2)\,Z\,\langle r^{-3}\rangle$, su logaritmo es
$\text{const} + \ln Z + \ln\langle r^{-3}\rangle$. El ajuste es lineal en $y$, así que el exponente
de $\zeta_{\text{desnudo}}$ es exactamente uno más que el de $\langle r^{-3}\rangle$: 1.690 frente a
0.690.
\item \textbf{El reparto del crecimiento.} De ese exponente 1.69, una unidad viene del factor $Z$
explícito y 0.69 del momento inverso. Es decir, $1/1.69 \approx 0.59$: el 59\,\% del crecimiento logarítmico
procede de la carga nuclear y el 41\,\% restante de $\langle r^{-3}\rangle$. Esa es la frase
\emph{algo más de la mitad} de la Sección~\ref{sec:regimen_acoplamiento}.
\item \textbf{El ajuste es descriptivo.} Son cuatro puntos y dos parámetros. Los exponentes locales
varían bastante (1.42, 2.14 y 1.86 para $\zeta_{\text{calc}}$), y el $R^2$ de
$\langle r^{-3}\rangle$ es solo 0.93. El exponente global resume la tendencia; no es una ley que
el modelo prediga. Como referencia, el ion hidrogenoide crece como $Z^4$ con $n$ fijo.
\end{enumerate}

%======================================================================
\section{Por qué la energía se desvía a segundo orden}
\label{det:segundo_orden}

La Sección~\ref{sec:parametros_fina} dice que la energía, por ser estacionaria, se desvía a segundo
orden en el error del orbital, mientras que $T$ y $\langle r^{-3}\rangle$ lo hacen a primer orden.

Sea $\psi_0$ el punto estacionario del funcional de energía y $\psi = \psi_0 + \varepsilon\,\phi$ una
función cercana, con $\varepsilon$ pequeño. Desarrollando,
\begin{equation}
  E[\psi_0 + \varepsilon\phi] = E[\psi_0]
  + \varepsilon\,\underbrace{\delta E[\psi_0;\phi]}_{=\,0}
  + \mathcal{O}(\varepsilon^2),
\end{equation}
porque la primera variación se anula en un punto estacionario. Para cualquier otro operador $\hat{A}$
que no define el funcional no hay tal cancelación:
\begin{equation}
  \ev{\hat{A}}{\psi} = \ev{\hat{A}}{\psi_0} + 2\varepsilon\,\Re\mel{\psi_0}{\hat{A}}{\phi}
  + \mathcal{O}(\varepsilon^2).
\end{equation}
El término lineal solo se anula por casualidad, por ejemplo si $\phi$ es ortogonal a
$\hat{A}\psi_0$.

En la tesis, $\varepsilon$ es la desviación que deja la ortogonalización de Gram-Schmidt
(Sección~\ref{sec:gram-schmidt}). La energía la siente al cuadrado y $T$ y $\langle r^{-3}\rangle$
linealmente. Es la jerarquía de las Tablas~\ref{tab:resultados_energia_global} y
\ref{tab:parametros_radiales}: en el silicio la energía difiere de la referencia en $10^{-6}$ Hartree
y la energía cinética en $10^{-3}$. Por la misma razón el cociente del virial, que contiene $T$,
se desvía de 2 más que la energía de su valor de referencia.

%======================================================================
\section{Un análogo clásico del paso de \texorpdfstring{$LS$}{LS} a \texorpdfstring{$jj$}{jj}}
\label{det:clasico}

La Sección~\ref{sec:acoplamiento_ls} dice que $\mathbf{L}$ y $\mathbf{S}$ \emph{dejan de ser
constantes de movimiento estrictas}. El mismo fenómeno aparece en mecánica clásica con cuatro
momentos angulares $\mathbf{l}_1, \mathbf{l}_2, \mathbf{s}_1, \mathbf{s}_2$ que obedecen el
corchete de Poisson del momento angular, $\{a_i, a_j\} = \epsilon_{ijk}\,a_k$ para cada vector, y
cero entre vectores distintos, con un Hamiltoniano modelo
\begin{equation}
  H = A\,\mathbf{l}_1\cdot\mathbf{l}_2 + B\,\mathbf{s}_1\cdot\mathbf{s}_2
    + \lambda\,(\mathbf{l}_1\cdot\mathbf{s}_1 + \mathbf{l}_2\cdot\mathbf{s}_2).
\end{equation}
Los dos primeros términos imitan la repulsión electrostática, que acopla las órbitas entre sí y,
por el intercambio, los espines entre sí; el tercero es el espín-órbita.

\begin{itemize}
\item \textbf{Siempre conservadas.} Las magnitudes $|\mathbf{l}_i|$ y $|\mathbf{s}_i|$, porque son
los invariantes (casimires) del corchete: el espacio de fases de cada vector es una esfera de radio
fijo, con su forma simpléctica de área. Y el momento total $\mathbf{J}$, por Noether, porque $H$ no
cambia al girar los cuatro vectores a la vez.
\item \textbf{$\lambda = 0$ (límite $LS$).} $H$ depende solo de $\mathbf{l}_1\cdot\mathbf{l}_2$ y
$\mathbf{s}_1\cdot\mathbf{s}_2$, así que tampoco cambia al girar las órbitas y los espines
\emph{por separado}. La simetría es mayor, $SO(3)_L\times SO(3)_S$, y Noether da dos cantidades
conservadas: $\mathbf{L}$ y $\mathbf{S}$.
\item \textbf{$A = B = 0$ (límite $jj$).} Ahora la simetría es girar $(\mathbf{l}_1,\mathbf{s}_1)$ y
$(\mathbf{l}_2,\mathbf{s}_2)$ por separado, $SO(3)_{j_1}\times SO(3)_{j_2}$, y se conservan
$\mathbf{j}_1$ y $\mathbf{j}_2$.
\item \textbf{Régimen intermedio.} Ninguna de las dos simetrías mayores sobrevive; solo queda la
rotación conjunta y solo $\mathbf{J}$ se conserva. Si $\lambda \ll A, B$, $\mathbf{L}$ y $\mathbf{S}$
precesan lentamente alrededor de $\mathbf{J}$, y sus magnitudes, que ya no son constantes, cambian
poco a lo largo de una precesión: son invariantes aproximados. Esa es la imagen clásica de las
etiquetas $LS$ que valen de manera aproximada.
\end{itemize}

El paso de $LS$ a $jj$ es entonces un cambio de cuál simetría mayor rompe la perturbación: el
espín-órbita rompe $SO(3)_L\times SO(3)_S$ hasta $SO(3)_J$, y la repulsión rompe
$SO(3)_{j_1}\times SO(3)_{j_2}$ hasta el mismo grupo. Las constantes de movimiento de cada límite
son los generadores de su grupo de simetría.

%======================================================================
\input{chapters/_borradores/polarizacion-capas-internas}

%======================================================================
\input{chapters/_borradores/estado-espurio-ci}

\printbibliography[heading=bibintoc]
\end{document}
